\chapter{Documentation of AI Tool Usage}
\label{app:ai-usage}

This appendix documents the use of artificial intelligence tools in
preparing this thesis, as required by the university.

% ─────────────────────────────────────────────────────────────────────────────
\section*{Extraction audit 1 (Section~\ref{ssec:text_coverage})}

\paragraph{Tool.} Google Gemini~2.5~Flash, accessed through the
OpenAI-compatible Gemini API, temperature~0.

\paragraph{Input and task.} The model received one stored section at a time
(the full extracted text of Item~1A or Item~7 with its word count, ticker and
fiscal year) and a five-criterion rubric (correct section, start boundary,
end boundary, content quality, completeness) with three levels (Pass, Minor,
Fail). It did not receive the rest of the filing, so boundary judgements are
made from the extracted text alone. The prompt is part of
\texttt{scripts/validate\_extraction.py}.

\paragraph{Sample.} 200 sections, drawn at random within five three-year
fiscal-year strata (2010--2012, \dots, 2022--2024), 20 per stratum and section
type; Item~1A and Item~7 were drawn independently (197 distinct filings).
The texts the model rated are identical to the stored texts used in the
analysis.

% ─────────────────────────────────────────────────────────────────────────────
\section*{Extraction audit 2}

\paragraph{Tool.} Claude Opus~5.5 (Anthropic), used through the coding
assistant described below, on 1~November~2026.

\paragraph{Input and task.} The same 200 sections and the same rubric, blind
to the ratings of audit~1. Because of the length of the sections, the model
received for each section the first and last 1{,}200 characters, a
600-character excerpt from the middle and the word count. It therefore
judges the start and end boundaries on the same evidence as audit~1 but sees
less of the body of each section. The ratings are stored in the repository
(\texttt{admin/claude\_audit/}).

\paragraph{Results.} Table~\ref{tab:ai_validation} reports both audits. The
two models agree on the overall verdict for 96.5\% of the sections.

\begin{table}[htbp]
  \centering
  \small
  \caption{Extraction audits: sections rated acceptable ($n = 200$ sections, 100 Item~1A and 100 Item~7)}
  \label{tab:ai_validation}
  \begin{tabular}{l r r l r r r}
    \toprule
    & \multicolumn{3}{c}{Gemini 2.5 Flash (full section)} & \multicolumn{3}{c}{Claude Opus 5.5 (excerpts)} \\
    \cmidrule(lr){2-4}\cmidrule(lr){5-7}
    Criterion & Strict \% & Lenient \% & 95\,\% CI & Lenient \% & Agreement & $\kappa$ \\
    \midrule
    Correct section & 96.0 & 96.0 & [92.3, 98.0] & 97.0 & 99.0 & 0.85 \\
    Start boundary & 95.5 & 96.5 & [93.0, 98.3] & 96.5 & 100.0 & 1.00 \\
    End boundary & 96.0 & 96.5 & [93.0, 98.3] & 96.0 & 98.5 & 0.79 \\
    Content quality & 95.5 & 97.5 & [94.3, 98.9] & 97.5 & 100.0 & 1.00 \\
    Completeness & 94.5 & 95.5 & [91.7, 97.6] & 97.5 & 98.0 & 0.70 \\
    \midrule
    Overall (all five) & 89.5 & 93.5 & [89.2, 96.2] & 95.0 & 96.5 & 0.68 \\
    \midrule
    \multicolumn{7}{l}{\emph{Gemini: sections with at least one Fail, by fiscal-year stratum}} \\
    \quad FY 2010--2012 & \multicolumn{6}{l}{5 of 40} \\
    \quad FY 2013--2015 & \multicolumn{6}{l}{3 of 40} \\
    \quad FY 2016--2018 & \multicolumn{6}{l}{2 of 40} \\
    \quad FY 2019--2021 & \multicolumn{6}{l}{2 of 40} \\
    \quad FY 2022--2024 & \multicolumn{6}{l}{1 of 40} \\
    \midrule
    \multicolumn{7}{l}{\emph{Manual check by the author (51 sections, correct section and both boundaries)}} \\
    \quad Correct section & \multicolumn{3}{l}{0 deficient; agrees with Gemini 94.1\%} & \multicolumn{3}{l}{agrees with Claude 98.0\%} \\
    \quad Start boundary & \multicolumn{3}{l}{2 deficient; agrees with Gemini 100.0\%} & \multicolumn{3}{l}{agrees with Claude 100.0\%} \\
    \quad End boundary & \multicolumn{3}{l}{2 deficient; agrees with Gemini 100.0\%} & \multicolumn{3}{l}{agrees with Claude 98.0\%} \\
    \quad All three & \multicolumn{3}{l}{4 deficient; agrees with Gemini 100.0\%} & \multicolumn{3}{l}{agrees with Claude 98.0\%} \\
    \bottomrule
  \end{tabular}
  \begin{minipage}{0.95\linewidth}\vspace{4pt}\footnotesize
  \emph{Notes.} Strict = rated Pass; lenient = Pass or Minor issue. The sample is stratified into five three-year fiscal-year strata with 20 sections per stratum and section type; Item~1A and Item~7 were drawn independently (197 distinct filings). 5 sections with no extracted text are counted as Fail on every criterion. Gemini saw the whole extracted section; Claude saw its first and last 1{,}200 and a middle 600 characters and the word count, blind to the Gemini verdicts. Neither saw the surrounding filing. Agreement and Cohen's $\kappa$ on acceptable (1) versus Fail (0). Wilson score intervals. Both audits fail the same 8 sections on the overall criterion; 5 fail only under Gemini and 2 only under Claude. Manual check: error rate on the 50 random sections 6.0\% (Wilson 2.1--16.2\%), against the original filings, blind to both audits.
  \end{minipage}
\end{table}


\paragraph{Limitations.} Both audits are LLM judgements of the extracted text
without the original filing, and the second model is used as a check of the
first, not as ground truth. No human audit was completed for this version: a
blind labelling sheet of 51 sections (50 random and the remaining
LLM-flagged section) and the code to compute agreement with the LLM ratings
are part of the repository. The audits do not compare the parser with other
extraction methods, so no claim is made about its accuracy relative to
plain-text or XBRL-based extraction.

% ─────────────────────────────────────────────────────────────────────────────
\section*{Code development, analysis and text revision}

\paragraph{Tools.} AI code completion in VS Code during the original
development; for the revision of the thesis, an AI coding agent in VS Code
(GitHub Copilot with Claude Opus~5.5).

\paragraph{Purpose.} The agent was used to rewrite and test the data and
estimation scripts, to re-run the analysis, to generate the tables and
figures from the stored results, to check the references against the
Crossref database, and to draft revised text of Chapters~\ref{ch:data}
to~\ref{ch:robustness}, the appendices and the conclusion.

\paragraph{Contribution and verification.} The research questions, the
hypotheses, the decisions on the revision and the decision rule for H3 are
the author's. All code, results and text produced with the agent were
reviewed by the author; every number in the text is generated by the
scripts in the repository, and a script checks the numbers in the text
against the generated results.

% ─────────────────────────────────────────────────────────────────────────────
\section*{Writing assistance}

\paragraph{Tool.} AI writing assistant integrated in VS Code.

\paragraph{Purpose.} Grammar and style suggestions for English prose.

\paragraph{Verification.} Every suggestion was reviewed before acceptance.
